\documentclass[10pt,a4paper]{amsart}
\usepackage[margin=19mm]{geometry}
\usepackage{amsmath,amssymb,mathtools}
\usepackage[scr=rsfs,cal=euler]{mathalfa}
\usepackage{xfrac}
\usepackage[unicode,hidelinks,bookmarks=false]{hyperref}
\newcommand{\ie}{i.e.\ }
\newcommand{\cf}{cf.\ }
\newcommand{\resp}{respectively\ }
\newcommand{\Subsection}{\subsection}
\providecommand{\nobreakdash}{\nobreak\hbox{-}}
\setlength{\emergencystretch}{2em}
\multlinegap0pt
\usepackage{amssymb}
\usepackage{tikz}
\newtheorem{prop}{Proposition}
\newtheorem{lem}{Lemma}
\newcommand{\B}{\mathcal{B}}
\newcommand{\C}{\mbox{${\mathbb C}$}}
\newcommand{\Wn}{{W}_n}
\newcommand{\cW}{{ W}}
\title{Orbits on a product of two flags and a line and the Bruhat order, II}
\author{Mark Colarusso, Sam Evens}
\date{Source: arXiv:2606.30478v1 (CC BY 4.0)}
\begin{document}
\maketitle

\noindent\emph{Source and licence.} Orbits on a product of two flags and a line and the Bruhat order, II. Version arXiv:2606.30478v1 (CC BY 4.0), distributed by arXiv under the Creative Commons Attribution 4.0 International licence, http://creativecommons.org/licenses/by/4.0/, as recorded in the arXiv licence metadata for this exact version and shown on the abstract page https://arxiv.org/abs/2606.30478v1. Full licence notice: \texttt{licenses/arxiv-2606.30478-CC-BY-4.0.txt}.\par\smallskip

\noindent\textit{Editorial scope.} Counted unit: the article's lem l:compare (source0.tex lines 2531--2537). The article's complete original proof of it is delivered with it (source0.tex lines 2538--2585, one \texttt{proof} environment of 48 lines). No mathematics is written, repaired or re-derived: the statement and the proof are the article's own lines, in the article's own notation, and no retained line is substituted or re-flowed.  Retained uncounted as the source-local context the proof needs: lines 2429--2431; lines 2444--2453; lines 2484--2486; lines 2488--2491; lines 2493--2495.  Nothing was omitted from any retained span, so the package carries no omission record.  No retained line contains a citation, so the wrapper carries no bibliography and no bibliography entry.  No retained line contains a figure environment, an image-inclusion command or drawing code, so no image is delivered and the package declares no asset dependency.  Editorial formatting only, in addition: an AMS \texttt{amsart} preamble that restates the article's own declarations and macros used by the retained lines, namely the environments \texttt{prop}, \texttt{lem} on separate counters, together with the standard packages those lines need.  The retained environments print in this document as Proposition 1, Lemma 1 in order of appearance, whereas the article prints the same items with the numbering of its own full document.  The complete native source of the article as distributed in the arXiv e-print for this exact version is bundled unchanged as \texttt{sources/204023/source0.tex}.
\begin{equation}\label{eq:orbitsembedding}
G\backslash (\B_{n}\times\B_{n}\times \mathbb{P}^{n-1})\hookrightarrow G_{n+1}\backslash (\B_{n+1}\times\B_{n+1}\times \mathbb{P}^{n}).
\end{equation}

\begin{prop}\label{p:marked}
Let $(w,\Delta)\in\widehat{\cW}_{n}$ and consider the $n+1$-cycle $\sigma=(n+1, n, n-1,\dots, 2, 1)$ of $\cW_{n+1}.$    Let $\phi: \widehat{\cW}_{n}\hookrightarrow \widehat{\cW}_{n+1}$ be the embedding on decorated permutations corresponding to the orbit embedding in Equation (\ref{eq:orbitsembedding}).  Then 
\begin{equation}\label{eq:permutation}
\phi((w,\Delta))=(w^{-1}\sigma, [w^{-1}(\Delta):n+1]).
\end{equation}
In particular, 
\begin{equation}\label{eq:opdecorated}
\mathcal{O}_{(w,\Delta)}^{op}=\mathcal{O}_{\phi(w,\Delta)}.
\end{equation}
\end{prop}

\begin{equation}\label{eq:firststat}
\mbox{For } p,\, q\in\{0,\dots, n\}\mbox{ define } r_{p,q}(w):=\dim(\mathcal{E}_{p}\cap w(\mathcal{E}_{q}))=|\{1,\dots, p\}\cap \{w(1),\dots, w(q)\}|.
\end{equation}    

\begin{equation}\label{eq:littledelta}
\mbox{ For } p,\,q\in\{0,1,\dots, n\}, \mbox{ define } \delta_{p, q} (w,\Delta):=\left\{\begin{array}{ccc} 1 & \mbox { if for each } \ell\in\Delta, & \ell\leq p \mbox{ or } w^{-1}(\ell)\leq q \\
0 &\mbox{ otherwise } & \end{array}\right\}_{.}
\end{equation}

\begin{equation}\label{eq:secondstat}
\overline{r}_{p,q}(w,\Delta):=r_{p,q}(w)+\delta_{p,q}(w,\Delta)\mbox{ for any } p,\, q\in\{0,1,\dots, n\}.
\end{equation}

\begin{lem}\label{l:compare}
Let $p\in\{0,\dots, n\}$ and $q\in\{1,\dots, n+1\}$.  %(MARK: CHECK THESE BOUNDS).
\begin{equation}\label{eq:shift}
\overline{r}_{p,q}(\phi(w,\Delta))=\overline{r}_{q-1,p} (w,\Delta).
\end{equation}
%where the statistic $\overline{r}_{p,q}$ is defined in Equation (\ref{eq:secondstat}).
\end{lem}

\begin{proof}
% Assume that $p,q \in \{ 0, \dots, n\}.$ 
 For $1\le k \le n+1,$ let ${\mathcal{E}}^{n+1}_{k}$ be the span of $e_1, \dots, e_k$ in $\C^{n+1}.$  Recall from Proposition \ref{p:marked} the element $\sigma^{-1} w \in {\cW}_{n+1}.$  We claim that
\begin{equation}\label{eq:rpqshift}
r_{p,q}(w)=r_{p+1, q}(\sigma^{-1} w)\mbox{ for } p, \, q\in\{0,\dots, n\},
\end{equation}
% where $r_{p,q}(w)$ is defined in (\ref{eq:firststat}) and
where as above, $r_{p+1, q}(\sigma^{-1}w)=\dim(\mathcal{E}^{n+1}_{p+1}\cap \sigma^{-1}w(\mathcal{E}^{n+1}_{q})).$  Since $w\in \Wn$, for any index $\ell\in\{1,\dots, q\}$, 
$1\leq w(\ell)\leq n$, so that $\sigma^{-1}w(\ell)=w(\ell)+1$.  
Therefore, for $p\geq 1$, we have a bijection between the sets $\{1,\dots, p\}\cap \{w(1),\dots, w(q)\}$ and $\{1,\dots, p+1\}\cap \sigma^{-1}w (\{1,\dots, q\})$ given by $x\mapsto x+1,$ which verifies (\ref{eq:rpqshift}) when $p \ge 1$ and $q \ge 1.$  In case $p=0,$ then $r_{0,q}(w)=0$, and since $\sigma^{-1}w(\ell)\geq 2$ for all $\ell\in\{1,\dots, n\}$, $r_{1,q}(\sigma^{-1} w)=0.$  For $q=0$, both sides of (\ref{eq:rpqshift}) are zero and the claim is trivial, so that (\ref{eq:rpqshift}) is verified in all cases.   

Suppose now that $p\in\{0,\dots, n\}$ and $q\in\{1,\dots, n+1\}.$  We assert that 
\begin{equation} \label{eq:deltashift}
\delta_{p,q}(\phi(w,\Delta))=\delta_{p,q}(w^{-1}\sigma, [w^{-1}(\Delta): n+1])=\delta_{q-1, p}(w,\Delta).
\end{equation}
%MARK:YOU NEED TO BE CAREFUL WITH THE LIMITS HERE.
%Suppose that %$\Delta=\{j_{1}<\dots<j_{k}\}\subset\{1,\dots, n\}.$
The first equality is from Proposition \ref{p:marked}, and for the second equality, by Equation (\ref{eq:littledelta}), we have 
\begin{equation}\label{eq:firstdelta}
\delta_{p,q}(w^{-1}\sigma, [w^{-1}(\Delta): n+1])=\left\{\begin{array}{ccc} 1 & \mbox { if for each } \ell\in\Delta, & w^{-1}(\ell)\leq p \mbox{ or } \sigma^{-1}w(w^{-1}(\ell))\leq q \\
&\;\;\;\;\;\;\;\;\;\;\;\;\;\;\;\;\;\;\;  & \mbox{ and } n+1\leq p\mbox{ or } \sigma^{-1}w(n+1)\leq q\\
0 &\mbox{ otherwise } & \end{array}\right\}_{.}
\end{equation}
Because $\Delta\subset\{1,\dots, n\}$, $w\in\Wn,$ and $\sigma^{-1}(n+1)=1,$  (\ref{eq:firstdelta}) becomes
\begin{equation}\label{eq:seconddelta}
\delta_{p,q}(w^{-1}\sigma, [w^{-1}(\Delta): n+1])=\left\{\begin{array}{ccc} 1 & \mbox { if for each } \ell\in\Delta, & w^{-1}(\ell)\leq p \mbox{ or } \ell\leq q-1 \\
&\;\;\;\;\;\;\;\;\;\;\;\;\;\;\;\;\;\;\;   & \mbox{ and } 1\leq q\;\;\;\;\;\;\;\;\;\;\;\;\;\;\;\;\;\;\;\;\;\;\;\;\;\\
0 &\mbox{ otherwise } & \end{array}\right\}_{.}
\end{equation}
Since $q\geq 1$, the second condition is always satisfied and the second equality in (\ref{eq:deltashift}) now follows from the definition of $\delta_{p,q}$ in (\ref{eq:littledelta}). 
Note that for any $y\in \cW_{n+1}$ and any indices $p,\, q\in \{0,\dots, n+1\}$, it follows from the definition of $r_{p,q}(y)$ in (\ref{eq:firststat}) 
%and (\ref{eq:secondstat})
 that 
\begin{equation}\label{eq:rpqinv}
r_{p,q}(y)=r_{q,p}(y^{-1}).
\end{equation}  
% Recall the definition of the statistic $\overline{r}_{p,q}$ in (\ref{eq:secondstat}).  
For all $p, q$ as above,  
\begin{equation*}
\begin{split}
 \overline{r}_{p,q}(w^{-1}\sigma, [w^{-1}(\Delta): n+1])&=r_{p,q}(w^{-1}\sigma)+\delta_{p,q}(w^{-1}\sigma,  [w^{-1}(\Delta): n+1])\\
&=r_{q,p}(\sigma^{-1}w)+\delta_{p,q}(w^{-1}\sigma,  [w^{-1}(\Delta): n+1]) \mbox{ by } (\ref{eq:rpqinv})\\
&=r_{q-1, p}(w)+\delta_{q-1,p}(w,\Delta)\mbox{ by } (\ref{eq:rpqshift}) \mbox{ and }(\ref{eq:deltashift})\\
&=\overline{r}_{q-1, p}(w, \Delta),
\end{split}
\end{equation*}
which by Proposition \ref{p:marked}, establishes Equation (\ref{eq:shift}) and the Lemma.
\end{proof}

\end{document}
